\scanpage{133}
\mspage{4}
% archive sheet number: 132
Cas commun à 2\textsuperscript{o} et 3\textsuperscript{o} :

correspondance
\diagram{\begin{tikzcd}[ampersand replacement=\&,column sep=large,row sep=large]
\Gamma \arrow[r,hook,"\unsure{f'|g'}"] \arrow[d,"f"'] \& E' \\
E \&
\end{tikzcd},}
telle que $\Gamma\unsure{\longrightarrow|{-}}E'$ soit \deleted{\unsure{c}}injectif\inserted{!!}. Dans ce cas qui \unsure{généralise} 2\textsuperscript{o} et 3\textsuperscript{o} (en fait, une \unsure{composition} des deux) on a encore
\[
A\longrightarrow A'\cup\{\varnothing\}.
\]

% margin note written slanted, rising along the left edge and overlapping the main text.
\margin{Le test si on est dans le cas 2\textsuperscript{o}, est $\deleted{\unsure{f_*|\alpha_*}}\unsure{\alpha_*|f_*}(1)=1$. Le test pour le cas 3\textsuperscript{o} est quoi ?? Pas l'injectivité sur les supports (C'est le fait que $f$ soit injectif, \unsure{sans} être néc. iso) ni la \unsure{fidélité de} la conservation de $|\unsure{\circ|\bullet}|$ (\unsure{$=$} \unsure{univoque}), $\unsure{C|\mathcal{C}}$ !! alors quoi ? Le fait que $\beta_*$ \unsure{commute à} $\alpha_*$. (En plus de l'injectivité de l'application sur les supports et de la conservation de $|\circ|$, ou de $S\cap T$, \unsure{$S\cap C T$|$S\cap\complement T$}\ldots\unsure{$c$|$C$}).}

Il ne semble pas y avoir d'application directe de figures, pour une application $g:E\longrightarrow E'$ qui ne soit injective.

Dans le cas d'un morphisme d'ateliers quelconque, il faut voir si se \unsure{construit|induit} tjs une ``correspondance'' sur les \unsure{spatios} associés --- d'où un $\unsure{\alpha_*|\alpha^*}$, $\unsure{\beta_*|\beta^*}$.

